% !TeX program = xelatex
\documentclass[UTF8,11pt,letterpaper,fontset=fandol]{ctexart}
% Standard TeX Live packages; no publisher-specific class or shell escape.
\setmainfont{Latin Modern Roman}
\usepackage[margin=1in]{geometry}
\usepackage{amsmath,amssymb}
\usepackage{graphicx}
\usepackage{booktabs,array,tabularx,longtable}
\usepackage{microtype}
\usepackage{enumitem}
\usepackage[font=small,labelfont=bf]{caption}
\usepackage[unicode,hidelinks]{hyperref}
\hypersetup{
  pdftitle={面向单风场多步功率预测的历史轨迹条件化天气修正},
  pdfsubject={单风场预测；CBH-SL 本站部署与 E2C 区域扩展；实验重审工作稿},
  pdfcreator={XeLaTeX}
}
\graphicspath{{figures/}}
\setlength{\emergencystretch}{2em}
\setlist[itemize]{leftmargin=*,itemsep=0.3em,topsep=0.4em}
\setcounter{topnumber}{2}
\renewcommand{\topfraction}{0.9}
\renewcommand{\textfraction}{0.1}
\renewcommand{\floatpagefraction}{0.75}
\DeclareMathOperator{\Concat}{Concat}
\DeclareMathOperator{\GELU}{GELU}
\DeclareMathOperator{\LN}{LN}
\DeclareMathOperator{\softplus}{softplus}
\DeclareMathOperator{\Readout}{Readout}
\DeclareMathOperator{\Decoder}{Decoder}
\DeclareMathOperator{\Encoder}{Encoder}
\DeclareMathOperator{\MLP}{MLP}
\newcommand{\R}{\mathbb{R}}
\newcommand{\trans}{\mathsf{T}}

\ctexset{section={format=\Large\bfseries},subsection={format=\large\bfseries}}
\renewcommand{\arraystretch}{1.15}

% Expandable file input for booktabs rows in alignment environments.
\makeatletter
\let\tableinput\@@input
\makeatother

\title{面向单风场多步功率预测的\\历史轨迹条件化天气修正}
\author{}
\date{}
\begin{document}
\maketitle
\begin{abstract}
单风场短期功率预测需要同时利用已观测运行历史与未来数值天气预报，并区分跨风场训练和部署时对其他风场数据的依赖。本文研究历史轨迹条件化的天气修正：先由历史预测器生成多步参考轨迹，再将参考水平、相邻差分和历史记忆均值显式提供给天气适配器；完整历史记忆独立支持修正解码，最终输出为参考与修正之和。在 56 个源风场上，预测任务为 24 个十五分钟历史步到未来 24 步。共享训练、仅使用本站在线输入的 CBH-SL，在 2023 验证与 2024 跨年复评中的平均容量因子 MAE 分别为 0.084140 和 0.081799，相对本批共享命名基线中验证集最佳的 Transformer 分别降低 4.09\% 和 3.42\%。共享/独立训练与本站/区域输入的四格比较表明，共享训练解释了较大的平均收益。区域历史版本的匹配条件消融支持显式轨迹信息包；非线性天气编码扩展 E2C 在 2024 达到 0.080703，但其结果不等同于本站部署版本。对 20 个未见风场的适配实验显示，源模型迁移相对从头训练具有样本效率优势，却未稳定优于预训练 Transformer。本文同时呈现按训练期描述量选取的代表风场、分场景误差及失败案例，以明确方法的适用范围。
\end{abstract}
\noindent\textbf{关键词：}风电功率预测；单风场部署；数值天气预报；轨迹条件；跨风场适配
\section{引言}
\label{sec:introduction}
单风场功率预测的输出是某个目标场站未来的发电轨迹，但这一输出定义并不限定模型必须只用该站的数据训练。跨风场共享参数可以扩大训练样本，而部署时仍可能只需目标风场的历史功率、本站天气预报与静态信息。与之不同，若模型在线读取区域聚合功率，则会增加对其他风场数据可用性的依赖。明确这两种信息来源，有助于判断预测收益能否转化为实际的单站部署能力。

历史功率与数值天气预报（NWP）提供互补的信息。前者反映已经实现的运行水平及短时变化，后者描述未来气象条件。在逐小时天气预测未来六小时、每十五分钟功率的任务中，插值可以对齐时间，却不能说明天气应如何改变已有的功率判断。本文关注的核心问题是：当历史模型已给出一条未来参考轨迹时，将这条轨迹显式传递给天气分支，是否有助于形成更准确的多步修正？

为此，本文先训练历史预测器，再冻结其参数。修正网络在每个提前量读取参考功率、相邻参考差分、提前量及历史记忆均值，联合对齐后的天气特征形成条件调整；解码器同时读取完整历史记忆。参考轨迹既是修正条件，又是最终输出的加性基准。这一组织方式将历史支持的未来判断显式暴露给天气分支。它没有增加独立观测，也不将历史模型解释为无天气影响的物理反事实；修正项允许同时承担天气相关预测和基准校准。

本文以共享训练、本站输入的 CBH-SL 为单站部署的主要实证对象，并用独立训练及区域历史输入构成四格比较。进一步在保留区域历史的版本上检验条件入口与非线性天气编码。这些实验回答不同问题：本地部署表现由 CBH-SL 支撑，模块归因由各自匹配的对照支撑，区域增强版本的优势不能直接转移为尚未评估的本站输入增强模型。类似地，源站误差更低也不意味着在新站上更容易适配，因此未见风场实验同时保留强预训练参照与负迁移实例。

本文的贡献概括如下：
\begin{itemize}
 \item \textbf{轨迹条件化的天气信息组织。}在保留完整历史解码记忆的同时，将历史模型的未来轨迹及其邻步差分显式用于天气修正。匹配条件对照识别这一信息包的收益，并将其与较小的摘要增量及单纯输出相加增量区分开。
 \item \textbf{训练共享与在线输入的分离验证。}以单风场预测为目标，通过共享/独立训练和本站/区域历史的四格设计，量化参数共享与在线区域输入各自对应的效果，为本站部署给出可复核的证据。普通参数共享本身不作为独立算法创新。
 \item \textbf{跨风场适配的收益及边界。}在不同适配训练预算下评价 20 个未见风场，并结合不依赖模型误差的代表站选择、完整场站统计和失败案例，揭示相对从头训练的收益与相对强预训练模型的不足。
\end{itemize}
主结论使用全部源站或最终目标站，不由少数典型站重新定义评价集合。天气编码的非线性、分时段表现与适配结果用于检验上述主张，不额外包装为未经独立验证的物理机制。

\section{方法}
\label{sec:method}
\subsection{单风场预测与信息边界}
设风场 $i$ 的功率为 $P_i(t)$，装机容量为 $C_i$，容量因子为 $p_i(t)=P_i(t)/C_i$。预测起点 $t$ 的历史与目标分别为
\begin{align}
 \mathbf H_{i,t}&=[p_i(t-(L-1)\Delta t),\ldots,p_i(t)],\\
 \mathbf y_{i,t}&=[p_i(t+\Delta t),\ldots,p_i(t+K\Delta t)],
\end{align}
其中 $L=K=24$、$\Delta t=0.25\,\mathrm h$，提前量 $\ell_k=k\Delta t$，最大提前量 $T=6\,\mathrm h$。24 个历史样本的首末时间间隔为 5.75 小时。$\mathbf W_{i,t}^{+}$ 表示起点可用的未来 NWP，$\boldsymbol\xi_i$ 为静态位置和容量，$\phi_t\in\{0,1,2,3\}$ 为小时内的起点相位。

本站输入设置只读取 $\mathbf H_{i,t}$、本站 NWP 和上述上下文。共享训练在多个源风场上优化同一组参数；独立训练则为每站单独估计参数。区域历史设置另外提供
\begin{equation}
 p_{\mathrm{agg}}(u)=\frac{\sum_{j\in\mathcal S}C_jp_j(u)}{\sum_{j\in\mathcal S}C_j},\qquad |\mathcal S|=56.
\end{equation}
本文以 S/I 表示共享/独立训练，以 L/C 表示本站/区域历史，得到 SL、IL、SC、IC。C 设置的区域历史仅进入历史预测器，修正侧的历史编码器仍只读取目标站历史。区域误差不是本文的主要评价目标。

\begin{figure}[tbp]
 \centering
 \includegraphics[width=\linewidth]{architecture_single_farm_v19.png}
 \caption{单风场轨迹条件化天气修正。历史预测器提供参考轨迹，独立历史编码器输出完整记忆及其均值摘要。条件适配器读取天气与轨迹条件；直接天气路径、目标上下文和完整记忆共同支持修正解码，最终为 $\hat{\mathbf y}=\mathbf b+\mathbf c$。下方区分已评估的 CBH-SL 本站输入版本和 E2C 区域历史扩展；后者额外向冻结预测器提供区域历史。图中曲线和特征为定性示意，不是实验数据；天气表示顺序为采样、时间对齐、编码。}
 \label{fig:overview}
\end{figure}

\subsection{历史支持的预测参考}
\label{sec:reference}
冻结历史预测器 $B_0$ 生成 $\mathbf b_i\in\R^K$。忽略起点下标，定义第二历史输入 $\mathbf G_i$：L 设置取 $\mathbf G_i=\mathbf H_i$，C 设置取 $\mathbf G_i=\mathbf H_{\mathrm{agg}}$。于是
\begin{equation}
 \mathbf b_i=B_0(\mathbf H_i,\mathbf G_i,\phi;\boldsymbol\xi_i).
 \label{eq:baseline}
\end{equation}
本站版本保留第二历史编码器的参数形状，但将其独立应用于本站序列，不通过其他站功率产生预测。固定的跨站静态归一化常数仍可使用，它们不构成在线区域功率输入。

参考模型的两条历史路径采用三个因果残差卷积块，膨胀率为 $1,2,4$。本站历史表示叠加静态嵌入，提前量与相位构成目标查询，带相对时间偏置的交叉注意力读取历史。该历史响应与第二路径的末位状态、目标嵌入融合，经未来序列上的三个残差卷积块生成更新 $\delta_{i,k}$。参考读出为
\begin{equation}
 b_{i,k}=\softplus\!\left[\softplus^{-1}\!\big(\max(p_i(t),10^{-4})\big)+\delta_{i,k}\right].
\end{equation}
未来卷积采用对称填充，属于联合多步预测，不是物理动力学递推。$B_0$ 不读取未来天气，但历史功率本身已包含过去气象和运行影响。

\subsection{历史记忆与显式轨迹条件}
修正侧历史编码器与参考模型独立。以历史功率和相对时间构成输入，生成 $d=64$ 维记忆：
\begin{align}
 \mathbf M_i&=\Encoder_H\left(\{\mathbf A_H[p_i(t+\tau_j);\tau_j/T]+\mathbf d_H\}_{j=1}^{L}\right),\\
 \tau_j&=(j-L)\Delta t,\qquad
 \mathbf m_i=L^{-1}\sum_{j=1}^{L}\mathbf M_{i,j}.
 \label{eq:summary}
\end{align}
$\Encoder_H$ 为两层、四头 Transformer，前馈宽度 128。全部历史在起点可用，因此历史自注意力无因果掩码。本文所报告 CBH 与 E2C 使用全部历史 token 的均值，未将末位 token 读出 SLAST 合并为当前实现。

参考差分为 $\Delta b_{i,1}=b_{i,1}-p_i(t)$，$\Delta b_{i,k}=b_{i,k}-b_{i,k-1}$（$k>1$）。它没有除以时间步长，是预测邻步变化而非未来真实爬坡。条件表示为
\begin{equation}
 \mathbf s_{i,k}=\LN\!\left(\GELU\left(\mathbf A_s[\mathbf m_i;b_{i,k};\Delta b_{i,k};\ell_k/T]+\mathbf d_s\right)\right)+\mathbf e_k.
 \label{eq:condition}
\end{equation}
$\mathbf e_k$ 为可学习的提前量偏置，跨风场及样本共享。完整记忆 $\mathbf M_i$ 另外直接进入解码器，因此中和额外摘要并不移除历史信息。

\subsection{天气对齐与天气编码}
以固定双线性权重从邻近四个网格采样 14 个天气变量，保留前七个逐小时节点的实际有效时间偏移 $\nu_{t,m}$。对每个提前量进行线性插值：
\begin{align}
 \mathbf z_{i,m}&=\sum_{g\in\mathcal N_i}\omega_{i,g}\mathbf W_{t,m}(g),\qquad \sum_g\omega_{i,g}=1,\\
 \bar{\mathbf z}_{i,k}&=(1-\alpha_{t,k})\mathbf z_{i,m}+\alpha_{t,k}\mathbf z_{i,m+1},\quad
 \alpha_{t,k}=\frac{\ell_k-\nu_{t,m}}{\nu_{t,m+1}-\nu_{t,m}}.
\end{align}
区间外保持最近端点值。这一操作对齐预测协变量，不将逐小时 NWP 解释为已恢复的十五分钟真实天气。

\textbf{本站主版本与线性对照。}CBH-SL、CBH-SC 和 FP32 对照 E0 采用线性天气路径。底层点式天气注意力每个提前量仅允许一个对齐 token，因而其 softmax 权重恒为一，值投影与输出投影可以合并为仿射映射。其有效形式为
\begin{equation}
 \mathbf w_{i,k}=\mathbf A_W\bar{\mathbf z}_{i,k}+\mathbf d_W.
\end{equation}
不将这一路径表述为基于查询的跨时间天气检索。

\textbf{区域非线性扩展。}E2C 替换天气输入投影为 $15\rightarrow64\rightarrow64\rightarrow64$ 的三层 MLP，中间使用 GELU，之后保留值及输出投影。输入为插值后天气和地形高程经可逆柱坐标变换、标准化后的 15 维表示。该变换包含地面变量与高程、低层相对地面风差及高度、两压力层间的风/温湿差与厚度；它没有增加原始天气与高程之外的信息。附录给出精确定义。E2N 使用原 14 个天气变量的同型非线性编码；E2L 以相同输入、层数和参数量移除中间激活，以检验非线性而非参数量的影响。E2C 是当前区域扩展的名称，不表示柱坐标已获得独立优势，也不代表本站输入 E2C 已完成评估。

目标上下文独立投影为 $\mathbf q_{i,k}=\mathbf A_q\boldsymbol\gamma_{i,k}+\mathbf d_q$，其中
\begin{equation}
 \boldsymbol\gamma_{i,k}=[\ell_k/T;\sin(2\pi\ell_k/24);\cos(2\pi\ell_k/24);\tilde\lambda_i;\tilde\eta_i;C_i/\bar C;\phi/3].
\end{equation}
正余弦编码提前量而非当地钟点；经纬度和容量采用固定源站统计归一化。

\subsection{条件天气修正与训练}
天气特征和轨迹条件经普通拼接适配器产生调整：
\begin{align}
 \mathbf a_{i,k}&=\mathbf A_2\GELU\!\left(\mathbf A_1[\mathbf w_{i,k};\mathbf s_{i,k}]+\mathbf d_1\right)+\mathbf d_2,\\
 \mathbf F_{i,k}&=\mathbf q_{i,k}+\mathbf w_{i,k}+\mathbf a_{i,k},\\
 \mathbf X_i&=\Decoder(\mathbf F_i,\mathbf M_i),\quad
 c_{i,k}=\Readout(\mathbf X_{i,k}),\quad
 \hat y_{i,k}=b_{i,k}+c_{i,k}.
 \label{eq:prediction}
\end{align}
适配器为 $128\rightarrow64\rightarrow64$，无显式乘法门控。两层四头解码器使用前置层归一化、128 维前馈层；标量读出为 $64\rightarrow32\rightarrow1$，中间使用 GELU。未来自注意力不加因果掩码，全部未来协变量在起点已可用；未来功率标签不作为输入。最终预测允许正负修正，不截断；特征调整也不解释为物理天气变化。

训练分两阶段：先以容量因子 MAE 训练历史预测器，选定 checkpoint 后将其冻结并固定为评估模式；再优化修正模型。共享训练目标为
\begin{equation}
 \mathcal L(\theta)=\frac1{|\mathcal T|NK}\sum_{t\in\mathcal T}\sum_{i=1}^{N}\sum_{k=1}^{K}
 |b_{i,t,k}+c_{i,t,k}-y_{i,t,k}|.
\end{equation}
独立训练则每次固定一个风场。修正读出和适配器末层零初始化，初始输出等于参考。第二阶段采用 AdamW，学习率 $2\times10^{-4}$、权重衰减 $10^{-4}$，余弦调度、梯度裁剪 1.0，最多 80 轮，验证 MAE 早停耐心值 10；dropout 为 0.05。不同实验系列的 BF16、FP16、FP32 设置分别报告，不将其跨系列差异视为单因素消融。模型未额外采用区域损失、爬坡损失或天气重建损失。

\section{实验结果与分析}
\label{sec:results}
\subsection{数据、协议与统计口径}
源站实验使用 56 个风场，2022 年训练、2023 年验证与选择模型，2024 年进行跨年复评。三个年份分别包含 31,190、32,550、28,782 个有效起报窗口。这里的全年协议指使用该年全部有效缓存，并非不间断的全年覆盖；2022 缓存截至北京时间 12 月 16 日，2024 年末有效窗口较少。2023 和 2024 已被多轮实验查看，2024 不称为未触碰盲测。

共享命名基线包含 15 个项目实现，均在全有效窗口上使用种子 42--51。H 表示历史输入家族，W 表示天气输入家族，FC 表示结合历史与未来协变量的家族；这些实现的输入组不同。CBH 的共享/独立对照采用 BF16，新四臂条件消融及输出匹配对照采用 FP16，七臂天气编码实验采用 FP32。主表为相同评价窗口下的\emph{完整系统比较}，各消融则只在其自身协议内解释。命名基线不等于已覆盖文献中所有官方实现或调参预算。旧独立命名基线仅使用按月份日期抽样的约 10\% pilot，训练/验证起点数为 3,156/3,349，且没有 2024 结果；它们独立列在附录，不能与全年结果混排。

令单模型在风场 $i$、种子 $s$ 的误差为 $e_{i,s,t,k}=\hat y_{i,s,t,k}-y_{i,t,k}$，本文主指标为
\begin{equation}
 \operatorname{MAE}=\frac1{NS}\sum_{i=1}^{N}\sum_{s=1}^{S}
 \frac1{|\mathcal T_i|K}\sum_{t\in\mathcal T_i}\sum_{k=1}^{K}|e_{i,s,t,k}|.
\end{equation}
先计算站内有效窗口与提前量平均，再等权平均种子和风场；报告的是单模型指标均值，不是集成预测误差。宏平均 RMSE 先对每个站、每个种子的 MSE 开平方，再平均；$R^2$ 同样在站内计算后平均。它们不同于先合并所有站 MSE 再开平方的统计量。轻量归档未保留的平方误差及 $R^2$ 不由 MAE 推算。

配对差分定义为 $\Delta(A,B)=\operatorname{MAE}(A)-\operatorname{MAE}(B)$，负值表示改善。除特别说明，场站区间先在站内平均种子，再以风场为重采样单位计算 95\% bootstrap 区间。天气扩展还采用共同七日块、种子与风场的交叉重采样，保持有效起点权重，并报告等周敏感性分析。区间描述对应样本和统计目标下的不确定性，不包含完整模型选择过程；场站相关性、跨年漂移及多重比较仍限制其解释。

\subsection{全场站主结果：本站部署与区域扩展}
\begin{table}[tbp]
 \centering\small
 \caption{56 个源风场的完整系统比较。报告十种子均值（Persistence 无训练种子）。2023 为验证，2024 为跨年复评；MAE/RMSE 越小越好，$R^2$ 越大越好。区域组额外读取区域历史，不能将两组差距归因于单一模块。破折号表示归档无法恢复该指标。}
 \label{tab:source-main}
 \setlength{\tabcolsep}{5pt}
 \begin{tabular}{lrrrr}
 \toprule
 & 2023 MAE & 2024 MAE & 2024 RMSE & 2024 $R^2$ \\
 \tableinput data/single_farm/source_main.tex 
 \bottomrule
 \end{tabular}
\end{table}
表~\ref{tab:source-main} 中，CBH-SL 的 2023/2024 MAE 为 0.084140/0.081799。相对验证集最佳共享命名基线 Transformer，分别降低 4.09\%/3.42\%；2024 在 49/56 个风场改善，场站配对差分区间为 $[-0.003651,-0.002095]$。因此，性能收益并不完全依赖部署时访问其他风场功率。它仍是两阶段条件修正系统与单阶段基线的比较，不能单凭该表将全部优势归因于条件入口。

区域扩展 E2C 的两年 MAE 为 0.082667/0.080703，相对共享 Transformer 降低 5.77\%/4.72\%，2024 在 50/56 站改善。但 E2C 同时具有区域历史输入、不同精度、天气表示和预训练过程，因此这两个百分比只描述系统表现。E2C 的模块增量应与同为区域输入、FP32 的 E0 比较。图~\ref{fig:horizon} 给出逐提前量表现；本站版本未保存可用于本次复核的最终 checkpoint，故不从区域模型曲线推断本站版本的场景行为。

\begin{figure}[tbp]
 \centering\includegraphics[width=\linewidth]{single_farm/03_forecast_horizons.pdf}
 \caption{源站逐提前量 CF MAE。命名基线为共享模型，E0/E2C 为区域历史版本；曲线覆盖各年全部有效窗口。该图用于展示系统随提前量的变化，其中 E2C 对 E0 的比较具有相同输入范围与精度。}
 \label{fig:horizon}
\end{figure}

\subsection{共享训练与区域历史的分离}
\begin{table}[tbp]
 \centering\small
 \caption{CBH 的共享/独立训练与本站/区域历史四格比较，十种子均值。所有单站预测均保留本站天气；C 的额外区域历史只进入参考模型。}
 \label{tab:sharing}
 \begin{tabular}{lllrr}\toprule
 配置 & 参数训练 & 在线功率输入 & 2023 MAE & 2024 MAE \\
 \midrule\tableinput data/single_farm/sharing.tex \bottomrule
 \end{tabular}
\end{table}
表~\ref{tab:sharing} 中，SL 相对 IL 的 2024 差分为 $-0.005199$，54/56 站改善；SC 相对 SL 的差分仅为 $-0.000291$。在这一协议中，跨站共享训练带来的平均差距显著大于额外区域历史的平均增量。IC 的均值也未优于 IL，说明更广的在线输入不会自动带来优势。四格比较支持将本站部署作为主要应用设置，不意味着所有时段或其他数据集上的区域历史都无价值。

\subsection{条件入口、输出形式与天气表示}
\begin{table}[tbp]
 \centering\small
 \caption{新四臂条件消融（区域历史，FP16，十种子重训）。所有臂均保留完整解码记忆、输出参考、提前量以及模块形状。CB 同时保留 $b$ 和 $\Delta b$。}
 \label{tab:condition}
 \begin{tabular}{lllrr}\toprule
 变体 & 额外历史均值摘要 & $b$ 与 $\Delta b$ & 2023 MAE & 2024 MAE \\
 \midrule\tableinput data/single_farm/condition.tex \bottomrule
 \end{tabular}
\end{table}
表~\ref{tab:condition} 的 CB 相对 CNEUTRAL 在 2024 降低 0.004419，支持轨迹与差分组成的条件信息包。CBH 相对 CH 降低 0.004593，说明在历史摘要之外显式提供这一信息包仍有价值。额外均值摘要的增量由 CBH 对 CB 识别，为 0.000260；它远小于轨迹信息包的增量。CH 对 CNEUTRAL 的均值差异很小，种子区间跨零，不将其提升为主要贡献。这里没有独立开关 $b$ 和 $\Delta b$ 的实验，因此不将两者作用分别量化。

独立的 matched direct/correction 对照使用相同 FP16 协议，2024 MAE 为 0.081631/0.081535。显式输出相加的差分为 $-0.000096$，场站区间 $[-0.000269,+0.000070]$ 包含零。主贡献应落在历史、轨迹与天气的信息组织上，单纯增加输出残差不足以解释主要收益。

\begin{table}[tbp]
 \centering\small
 \caption{天气编码七臂比较（区域历史，FP32，十种子）。E2(raw) 为原始天气加高程；E2C 是可逆坐标重参数化后输入 MLP。该组与表~\ref{tab:condition} 的精度不同，只作组内归因。}
 \label{tab:weather}
 \setlength{\tabcolsep}{5pt}
 \begin{tabular}{lllrr}\toprule
 变体 & 输入 & 编码器 & 2023 MAE & 2024 MAE \\
 \midrule\tableinput data/single_farm/weather.tex \bottomrule
 \end{tabular}
\end{table}
在表~\ref{tab:weather} 中，E2N 与 E2L 的输入、层数和参数量相同，差别是中间 GELU 激活。2023 差分为 $-0.001025$，支持非线性天气编码；2024 均值亦改善，但归档等周区间跨零。E2C 相对 E2(raw) 的两年增量仅约 0.000100/0.000101，不能据此确认柱坐标的独立稳健收益。高程与显式三分组也没有得到独立稳健支持，不将其列为主要创新。

E2C 对 E0 的 2024 有效起点加权差分为 $-0.000825$。仅重采样风场时的区间见附录表~\ref{tab:weights}。共同七日块、种子和风场交叉重采样得到 $[-0.001358,-0.000292]$。每周有效起点为 32--672，若改为各有观测周等权，点估计则为 $-0.000557$，区间为 $[-0.001178,+0.000168]$。两者对应不同统计目标；本文以与主 MAE 一致的有效起点权重为主，并保留等周敏感性，而非按区间是否跨零选择算法。

\subsection{典型风场与分场景结果}
典型站由训练期描述量选择：源站使用 2022 年、迁移目标站使用 2023 年，特征包括经纬度、对数容量、容量因子均值/标准差、低/高功率占比及过去一小时变化的 90\% 分位。各特征以中位数/IQR 缩放，地理、容量、功率分布、动态变化四组等权，通过确定性 PAM 选取源站六个、目标站四个中心站。此过程不读取模型误差，但属于本次重审确定的回顾性规则，不是原实验预注册；全部 56/20 站仍为主评价集合。原新增站队列的资格筛选曾利用 2024 覆盖率和跨年质量信息，与本次仅使用训练期描述量的代表站选择须分开说明。

\begin{table}[tbp]
 \centering\small
 \caption{六个训练期描述量代表源站的 2024 MAE，均为十种子均值。Transformer 与 CBH-SL 使用本站输入；E2C 额外使用区域历史。名称缩写仅用于排版。}
 \label{tab:representatives}
 \setlength{\tabcolsep}{4pt}
 \begin{tabular}{llrrrr}\toprule
 风场 ID & 名称 & MW & Transformer & CBH-SL & E2C \\
 \midrule\tableinput data/single_farm/source_representatives.tex \bottomrule
 \end{tabular}
\end{table}
表~\ref{tab:representatives} 涵盖大岗、马力、双龙、王府、到保和协和新发六站，覆盖容量、位置、运行分布和短时变化的差异。代表站用于呈现异质性，不按所属簇大小额外加权。完整逐站主结果见附录，位置和描述量分布见图~\ref{fig:selection}。

\begin{figure}[tbp]
 \centering\includegraphics[width=\linewidth]{single_farm/04_scenario_contrasts.pdf}
 \caption{区域天气扩展的探索性分场景比较，差分为 E2C 减 E0，负值表示改善。误差棒为场站配对 bootstrap 区间，未进行所有分组的多重比较校正。分组依据及其可用时点见正文；未来真实爬坡只用于事后诊断。}
 \label{fig:scenario}
\end{figure}
以下场景分析均比较区域 E2C 与 E0，不能替代本站主模型相对强基线的全年场景复评。2024 年，0--1 小时的平均改善为 0.000159，1--3 小时为 0.000620，3--6 小时为 0.001185，中远期增量更明显。按起报 UTC 月份分季，春季改善 0.001355，47/56 站改善；秋季均值改善的场站区间仍跨零。按北京时间起报时刻分为 00--06、06--12、12--18、18--24 四段，差分依次为 $-0.000442,-0.000714,-0.001018,-0.001131$；这与按未来目标发生小时分组不同。

起点功率和过去一小时变化均可在部署时观测。起点 CF $\geq0.7$ 的组内平均改善较大，但 2024 只有 55 站出现该状态，分母不是全部 56 站。另将未来六小时中任一真实一小时 CF 变化绝对值 $\geq0.2$ 的窗口定义为未来爬坡诊断组，E2C 对 E0 的差分为 $-0.001892$；非爬坡组为 $-0.000085$，区间跨零。该标签使用未来真值，不可用于在线路由，也不证明模型已准确识别爬坡的发生时间和幅度。

预测案例按固定规则取每站 2024 年 4 月北京时间 06:00 的首个满足普通或历史已观测爬坡条件的有效起点，不存在时按同条件顺延至全年首个起点。选择不使用未来真值或模型误差。图~\ref{fig:cases} 展示 seed42；E2C 阴影为十种子的 10--90\% 分位，表示训练随机性，不能当作预测置信区间。Transformer 案例由保存 checkpoint 在 CPU FP32 重新推理，其数值与归档 BF16 主指标不是同一精度的完整复评。附录保留事后附加的龙湖失败案例，不要求每幅图都显示本方法最好。

\begin{figure}[p]
 \centering\includegraphics[width=\linewidth]{single_farm/06_source_forecast_examples.pdf}
 \caption{六个代表源站的固定规则预测窗口。普通窗口与历史爬坡窗口均由已观测信息选取；实线为固定 seed42，阴影只表示 E2C 的种子间分位范围。E0/E2C 使用区域历史；这些曲线不代表缺少最终 checkpoint 的 CBH-SL。失败或偏差较大的窗口同样保留。}
 \label{fig:cases}
\end{figure}

\subsection{未见风场适配：数据效率与失败边界}
源站 56、开发站 10、最终目标站 20 互不重叠。目标站以 2023 年适配训练、2024 上半年选择 checkpoint、下半年评价。7/30/90 天预算各包含四个季节训练块，180 天和 full 各一个块。统一使用共同种子 42--46，先在每站内平均种子和训练块，再等权平均 20 站。Scratch 家族在开发站及上半年验证中选定，不在最终站下半年上逐预算挑选最优家族。CBH/E2C 的 T5 方案适配修正网络，保留冻结参考；其参考模型仍读取源 56 站区域历史，故不是纯本站输入迁移。

\begin{table}[tbp]
 \centering\small
 \caption{20 个未见风场在 2024 下半年的 MAE。预算仅为适配训练数据量，所有预算仍使用完整目标站上半年标签选择 checkpoint。E2C 为 FP32，历史迁移参照精度不同；不将微小差异作为精度完全匹配的归因结果。}
 \label{tab:transfer}
 \setlength{\tabcolsep}{4pt}
 \begin{tabular}{llrrrr}\toprule
 预算 & 开发集选定 scratch & Scratch & 预训练 TF & CBH-T5 & E2C-T5 \\
 \midrule\tableinput data/single_farm/transfer.tex \bottomrule
 \end{tabular}
\end{table}
E2C 对开发集选定 scratch 在五档预算下平均降低约 22.70\%、23.88\%、22.25\%、14.58\%、6.99\%，场站区间均支持改善（表~\ref{tab:transfer}）。这体现适配方案整体的样本效率，但两者的优化、预训练、冻结参考和输入路径不同，并非单一预训练开关的因果对照。预算只计训练标签，不能据 7 天结果声称总共仅需 7 天标注或严格冷启动。

强预训练 Transformer 改变了对迁移收益的解释：E2C 仅在 full 预算下具有更低的平均 MAE，而该差分区间仍跨零。30 天预算的 E2C 减 Transformer 为 $+0.001235$，95\% 场站区间 $[+0.000049,+0.002509]$。因此，源站上更好的 E2C 预测不构成其表示更可迁移的证据。

\begin{figure}[tbp]
 \centering\includegraphics[width=0.92\linewidth]{single_farm/05_transfer_learning_curve.pdf}
 \caption{最终 20 站适配训练预算曲线，统一共同五种子。天数仅计训练预算，上半年 checkpoint 选择标签未计入其中。图中较低误差代表更好；全数据均值的微小排序不等同于可靠迁移优势。}
 \label{fig:transfer}
\end{figure}
目标代表站为甄峰（49.5 MW）、水字平价（100 MW）、君字（300 MW）与同发龙源特许权（200 MW）。这些站并非均为有利案例：甄峰在 7 天预算下 E2C 的 MAE 为 0.121572，劣于预训练 Transformer 的 0.108034，也高于 scratch 的 0.116830；君字多数预算亦未优于 Transformer。图~\ref{fig:target-cases} 保留这些边界，避免将跨站平均收益误认为所有新站均受益。

\section{讨论与局限}
现有结果支持轨迹条件化的信息组织及共享训练下的本站部署，但证据来自不同协议：CBH-SL 支撑本站部署，条件与天气消融主要来自区域版本。本站 E2N/E2C 及其匹配消融尚未完成；当前归档也未恢复 SL 最终 checkpoint，完整逐样本复评仍受限。因此不能将这些结果拼接为已验证的本站增强模型。

后续需补齐本站增强结构与强基线的共同精度场景复评；迁移应计入 checkpoint 选择标签，或在开发站冻结训练步数。单一区域、缺测窗口、已查看年份与队列筛选限制外推；场站区间未覆盖全部空间依赖和模型选择，探索性分组也未统一校正。新确认应使用独立年份或站组。

\section{结论}
历史轨迹条件化天气修正在保留完整历史记忆的同时，显式提供待修正的未来功率判断。全 56 站结果支持共享训练、本站部署；匹配对照表明轨迹信息包的收益大于额外摘要及单纯输出相加的增量。区域非线性天气扩展进一步降低源站误差，但未见站适配未稳定超过强预训练 Transformer。区分训练来源、在线输入和标签预算，是解释这些收益及其适用范围的关键。

\clearpage
\appendix
\section{实现、协议和复核补充}
\subsection{输入和训练配置}
NWP 通道依次为海平面气压、10 m 风的 $u/v$ 分量、2 m 温度，以及 1000/925 hPa 上的位势、比湿、温度和 $u/v$ 风分量，共 14 维。插值保留实际预报有效时间，不能将未来天气观测代替起点可用预报。历史预测器的宽度为 64，两组历史卷积与未来卷积均为三层，参考训练 dropout 为 0.1；修正侧宽度 64、注意力头数 4、编码器和解码器各两层，dropout 为 0.05。历史参考预训练最多 40 轮，学习率 $3\times10^{-4}$；修正训练最多 80 轮。不同归档系列的实际 batch 组织及精度以各自配置为准，不能用一个默认值覆盖全部对照。

15 个共享命名基线每个家族/种子只有一套参数，复用项目中相应模型实现。它们不增加场站 ID、场站嵌入或区域历史；每轮使用全部 $(\mathrm{origin},\mathrm{station})$ 组合，共 1,746,640 个训练样本。采用 MAE、AdamW，学习率 $2\times10^{-4}$、权重衰减 $10^{-4}$、BF16、batch size 64、最多 80 轮、早停耐心值 10。2023 年 56 站等权 MAE 选择 checkpoint。本组没有证明各架构分别获得了穷尽的超参数搜索，因而主文将其称为本批实现的基线比较。

目标 T5 适配冻结历史参考，更新修正核心并学习逐提前量偏置；T5 不包含 T2/T4 的低秩适配器，归档名称中的 r16 不应解释为 T5 的有效低秩瓶颈。主表只比较已对齐的共同五种子，不将不同种子数量当作方法优势。所有预算使用目标站完整上半年验证标签，因此需要与总标签受限的少样本设置区分。

\subsection{E2C 坐标定义及归因控制}
令 $h$ 为地形高程，$z_p$ 为压力层位势，$g_0=9.80665$，$\theta_{1000}=T_{1000}$，$\theta_{925}=T_{925}(1000/925)^{0.2854}$。实际 15 维输入按以下顺序排列：
\begin{align}
 \mathbf v_{\mathrm{surface}}&=[p_{\mathrm{msl}},u_{10},v_{10},T_2,h],\\
 \mathbf v_{\mathrm{lower}}&=[u_{1000}-u_{10},v_{1000}-v_{10},\theta_{1000},q_{1000},z_{1000}/g_0-h],\\
 \mathbf v_{\mathrm{upper}}&=[u_{925}-u_{1000},v_{925}-v_{1000},\theta_{925}-\theta_{1000},
 q_{925}-q_{1000},(z_{925}-z_{1000})/g_0].
\end{align}
该变换可逆，各通道随后按训练统计标准化。E2C 将三组拼接后输入普通 MLP；E3 才显式采用 $5\rightarrow16$、$5\rightarrow24$、$5\rightarrow24$ 的分组支路。保留坐标和高程控制的目的在于避免把非线性编码的整体收益误归因于某个物理命名的输入。

\subsection{代表站、窗口与统计复核}
源站代表选择使用 2022 有效起点的功率描述量；目标站使用 2023。低/高容量因子描述量阈值为 $0.05/0.8$，与主文场景诊断的 $0.1/0.3/0.7$ 分箱不是同一用途。地理两维各按 $1/\sqrt2$、分布四维各按 $1/2$ 加权，使四个描述组在平方距离中等权。确定性 PAM 选择源六站和目标四站；簇大小仅描述覆盖范围，不改变论文指标的风场等权规则。

普通案例要求起点 $0.1\leq p_i(t)<0.7$ 且过去一小时变化绝对值不超过本站训练期 90\% 分位；历史爬坡案例要求该变化超过阈值。本次六个代表站与龙湖的全部 14 个窗口均在 2024 年 4 月找到，无需全年回退。病例展示日期由规则决定，未来真值仅在选定之后用于绘制预测效果。

本次结果整理核对了 40 个 E0/E2C 预测归档，重算 MAE 与记录的最大绝对差约 $5.4\times10^{-10}$；448 行季节结果的最大差约 $9.9\times10^{-10}$。新条件消融的 CB 输入通过源码及真实 checkpoint 的前向 hook 核验，确认包含 $b$ 和 $\Delta b$。图表重汇总与少量固定窗口推理不构成新的模型训练。

\begin{table}[tbp]
 \centering\small
 \caption{2024 E2C 减 E0 的统计目标敏感性。不同权重对应不同估计量，不能把等周区间配到有效起点均值上。}
 \label{tab:weights}
 \begin{tabular}{lrr}\toprule
 统计目标与重采样单位 & 差分 & 95\% 区间 \\
 \midrule
 有效起点等权；风场 & $-0.000825$ & $[-0.001214,-0.000441]$ \\
 有效起点等权；种子/风场/共同七日块 & $-0.000825$ & $[-0.001358,-0.000292]$ \\
 有观测周等权；种子/风场/共同七日块 & $-0.000557$ & $[-0.001178,+0.000168]$ \\
 \bottomrule\end{tabular}
\end{table}
其他天气因子归档区间沿用等周重采样，正文只报告明确的权重解释，不静默将它们替换为有效起点区间。场景分组、完整种子统计、月份覆盖数和目标逐站结果保存在随稿机器可读数据中。

\section{补充图：代表性、系统比较与失败边界}
\begin{figure}[p]
 \centering\includegraphics[width=\linewidth]{single_farm/01_farm_selection.pdf}
 \caption{代表风场的地理与训练期运行描述量。源站六个和最终目标站四个代表均由描述量聚类选取，不读取模型误差。代表集合不替代全 56/20 站评价；原目标队列资格筛选使用过跨年质量信息。}
 \label{fig:selection}
\end{figure}
\begin{figure}[p]
 \centering\includegraphics[width=\linewidth]{single_farm/02_source_baseline_comparison.pdf}
 \caption{源站系统比较补充视图。本站输入和区域历史系统分别识别；单站部署主结论由 CBH-SL 支撑，E2C 曲线为区域扩展。所有数值来自已核验实验记录。}
 \label{fig:source-summary}
\end{figure}
\begin{figure}[p]
 \centering\includegraphics[width=\linewidth]{single_farm/07_failure_forecast_examples.pdf}
 \caption{龙湖（1043056）的事后附加失败诊断。它不是描述量聚类选出的六个代表站之一；窗口仍遵循同一普通/历史爬坡规则。误差较大的实例用于说明边界，不纳入有利案例筛选。阴影为训练种子分位范围，不是预测区间。}
 \label{fig:failure}
\end{figure}
\begin{figure}[p]
 \centering\includegraphics[width=\linewidth]{single_farm/08_target_medoid_learning_curves.pdf}
 \caption{四个训练期描述量代表目标站的适配预算曲线。甄峰与君字保留相对强预训练 Transformer 的劣势，说明全站平均收益不意味着每站或每个预算均受益。天数为训练标签预算，不包含上半年 checkpoint 选择标签。}
 \label{fig:target-cases}
\end{figure}
\clearpage
\section{全部源风场与独立 pilot 结果}
以下表格保留全部 56 站，避免只呈现代表站。Transformer/IL/SL 使用本站在线输入，E0/E2C 使用区域历史；不跨输入组解释单模块因果增量。完整机器可读文件还包含 2023 指标、全部共享命名家族与 SC/IC。
\small
\begin{longtable}{lrrrrr}
 \caption{全部 56 个源风场的 2024 MAE，十种子均值。}\label{tab:all56}\\
 \toprule 风场 ID & Transformer & CBH-IL & CBH-SL & E0 & E2C \\
 \midrule\endfirsthead
 \multicolumn{6}{l}{表~\ref{tab:all56}（续）}\\
 \toprule 风场 ID & Transformer & CBH-IL & CBH-SL & E0 & E2C \\
 \midrule\endhead
 \midrule\multicolumn{6}{r}{续下页}\\\endfoot
 \bottomrule\endlastfoot
 \tableinput data/single_farm/source_all56.tex 
\end{longtable}
\normalsize
\begin{table}[tbp]
 \centering\small
 \caption{独立命名基线的 pilot 结果，仅有 2023 验证。训练/验证为按月份日期抽样的 3,156/3,349 起点；该表与全年协议的训练量和评价窗口均不同，不能混排。每模型为 56 站、十种子，RMSE/$R^2$ 为站内计算后宏平均。}
 \label{tab:pilot}
 \begin{tabular}{lrrr}\toprule
 模型 & MAE & RMSE & $R^2$ \\
 \midrule\tableinput data/single_farm/pilot.tex \bottomrule
 \end{tabular}
\end{table}
\section{未列为主贡献的结果与复现范围}
早期 SLAST 摘要读出、高频天气恢复、生成式天气模块及低秩站点个性化没有构成本稿的主方法。大型摘要未相对同容量均值控制显示清晰收益；高频恢复候选未超过对应线性对照；低秩迁移也未形成相对容量控制的稳定优势。这些历史结果用于解释设计取舍，不能替换当前匹配的条件消融或作为未验证增强重新并入最终结构。

随稿数据保留主结果、四格比较、新 FP16 条件实验、FP32 天气七臂、目标共同种子比较、场景统计与完整源站明细。完整实验重审还包含各系列来源路径与哈希、统计脚本及待补实验配置。框架图使用 imagegen 的 CLI 路径请求 \texttt{gpt-image-2.5} 生成并按源码核对；数值结果图由实验数据绘制，未使用生成图像替代定量结果。论文源文件与编译输入均可独立打包。

\end{document}
